\documentclass[11pt]{article}

\usepackage[margin=1in]{geometry}
\usepackage{amsmath,amssymb,amsthm,mathtools}
\usepackage{booktabs,tabularx,array}
\usepackage{microtype}
\usepackage{enumitem}
\usepackage{xcolor}
\usepackage{xurl}
\usepackage[hidelinks]{hyperref}

% RELEASE GATE: named authors, affiliations, ORCIDs, and a contribution
% statement replace the project line before circulation.

\newcommand{\ii}{\mathrm{i}}
\newcommand{\id}{\mathbf{1}}
\newcommand{\C}{\mathbb{C}}
\newcommand{\R}{\mathbb{R}}
\newcommand{\norm}[1]{\left\lVert #1\right\rVert}
\definecolor{kernelcolor}{RGB}{0,92,74}
\newcommand{\Kernel}{\textcolor{kernelcolor}{\textsc{Kernel}}}

\newtheorem{theorem}{Theorem}
\newtheorem{corollary}{Corollary}
\newtheorem{proposition}{Proposition}

\title{Machine-Verified Fermion-Doubling Obstructions for Ordered $3+1$
Dirac Quantum Walks}
\author{Mark Schwab}
\date{July 11, 2026 --- letter draft v1 (overnight run)}

\begin{document}
\maketitle

\begin{abstract}
Whether a strictly local, exactly unitary, discrete-time walk can simulate
the $3+1$-dimensional Dirac equation without fermion doubling is a live
design problem, with its lattice origin in the Nielsen--Ninomiya theorem
\cite{NielsenNinomiya}.  For the ordered successive-axis (split-step) walk --- the
standard minimal architecture --- we settle the spectral side of that
problem exactly, with every statement checked by the Lean~4 kernel.  We
prove closed determinant factorizations valid at every momentum and mass
angle, $\det(U\mp\id)=4P_{0,\pi}$ with $P_{0,\pi}$ explicit polynomials in
the momentum cosines, and, for every mass angle with $0<|\cos\theta|<1$, a
complete crossing classification: zero- and $\pi$-quasienergy modes occur
exactly where all three momentum cosines vanish.  Independently of that
branch restriction, the massive walk admits no uniform Floquet gap at any
angle: body-center crossings at quasienergies
$0$ and $\pi$ persist for every mass angle.  Two structural no-go theorems
delimit the possible repairs: a degree-one nearest-neighbor factor admits no
stationary amplitude compatible with origin normalization, exact unitarity,
and an involutory unit-speed tangent; and every momentum-independent onsite
coin leaves all three even-parity corner aliases of the massless walk
intact.  As a corollary, the recently proposed stay-put doubler-free family
of Gupta and Short, which is exactly of the covered degree-one form,
necessarily has a non-involutory tangent generator: stay-put freedom is
purchased by giving up the unit-speed Dirac tangent.  These are boundary
theorems, not a doubler-free construction; they convert the search for a
clean strict $3+1$ regulator into a precise, machine-checkable algebra
problem and supply its regression tests.
\end{abstract}

\section{Setup: the ordered walk, and what ``local'' means here}
\label{sec:setup}

Two locality notions must not be conflated.  A \emph{strict} discrete-time
walk updates the state by one application of a finite-range unitary; a
lattice Hamiltonian generates $e^{-\ii tH}$, which is exactly unitary but
not, in general, a finite-range one-step update.  Every claim in this
letter is about the strict notion unless explicitly labeled Hamiltonian.

Let $\alpha_1,\alpha_2,\alpha_3,\beta\in M_4(\C)$ be the standard mutually
anticommuting Hermitian involutions and let the mass coin carry angle
$\theta$.  The ordered successive-axis walk --- the architecture used by
Mlodinow and Brun \cite{MlodinowBrun} and the minimal $3+1$ lift of the
split-step walk --- has one-step Bloch symbol
\begin{equation}\label{eq:walk}
 U(\mathbf q,\theta)
 =e^{-\ii q_x\alpha_1}e^{-\ii q_y\alpha_2}e^{-\ii q_z\alpha_3}
  e^{-\ii\theta\beta},
 \qquad \mathbf q\in[-\pi,\pi]^3 .
\end{equation}
On a finite periodic register the corresponding position-space update is
exactly unitary, has an explicit two-sided inverse, and reads only sites
reached by one conditional shift per axis; the plane-wave sectors reproduce
\eqref{eq:walk} exactly \cite{CompanionDerivation}.  The factor ordering is
not cosmetic: the spatial generators do not commute.

Doubling and its variants in walk and automaton discretizations have a
growing exact literature \cite{MlodinowBrunNoGo,GuptaShortQCA}; alternative
$3+1$ architectures include the tetrahedral walk \cite{NzonganiTetra} and
the principle-derived Weyl/Dirac automata \cite{DArianoPerinotti,FreeQCA}.  The present
contribution is orthogonal to construction: it fixes, exactly and
machine-checkably, what the minimal ordered architecture cannot do, and
which repairs are closed.

All results below are kernel-checked in Lean~4 over Mathlib
\cite{Lean4,Mathlib}; each headline declaration carries a build-enforced
axiom footprint equal to
$\{\mathtt{propext},\mathtt{Classical.choice},\mathtt{Quot.sound}\}$.  The
declaration table is in the appendix.  The kernel checks the statements;
which physical readings they support is argued, as everywhere, in prose.

\section{An exact all-zone spectral criterion}
\label{sec:criterion}

Write $x=\cos q_x$, $y=\cos q_y$, $w=\cos q_z$, $c=\cos\theta$, and define
\begin{align}
 S={}&4c^2x^2y^2w^2-2c^2x^2y^2-2c^2x^2w^2+c^2x^2
      -2c^2y^2w^2+c^2y^2+c^2w^2 \notag\\
   &{}-2x^2y^2w^2+x^2y^2+x^2w^2+y^2w^2,
 \label{eq:S}\\
 P_0&=S-2cxyw,
 \qquad
 P_\pi=S+2cxyw.
 \label{eq:P}
\end{align}

\begin{theorem}[All-zone determinant factorization]\label{thm:det}
For every real momentum and mass angle,
\begin{equation}\label{eq:detfact}
 \det\!\bigl(U(\mathbf q,\theta)-\id\bigr)=4P_0,
 \qquad
 \det\!\bigl(U(\mathbf q,\theta)+\id\bigr)=4P_\pi .
\end{equation}
Consequently $P_0=0$ and $P_\pi=0$ are, respectively, necessary and
sufficient for a zero- and a $\pi$-quasienergy mode.
\end{theorem}
\noindent\Kernel{}\quad The identities are proved from the explicit
$4\times4$ expansion and bridged definitionally to the live walk used in
the companion paper's continuum estimates.

Corner sampling is not a classification; \eqref{eq:detfact} is.  The zero
sets themselves admit a complete description on the massive principal
branch:

\begin{theorem}[Complete crossing classification]\label{thm:class}
For every mass angle with $0<|\cos\theta|<1$,
\begin{equation}
 P_0=0
 \iff
 P_\pi=0
 \iff
 x=y=w=0,
\end{equation}
that is, zero- and $\pi$-quasienergy crossings of the massive ordered walk
occur exactly where all three momentum cosines vanish
($q_j\in\{\pm\pi/2\}$).
\end{theorem}
\noindent\Kernel{}\quad The proof is a sharp sum-of-squares/AM--GM
decomposition of $S\mp2cxyw$ on $|x|,|y|,|w|\le1$.  The boundary angles
carry their own controls: at the quarter-turn angle $\cos\theta=0$
explicit extra zeros are exhibited, and the massless angles
$\cos\theta=\pm1$ reduce to the corner classification of
Corollary~\ref{cor:aliases}.

\section{No-gap and alias obstructions}
\label{sec:obstructions}

\begin{theorem}[No uniform Floquet gap]\label{thm:bodycenter}
At body-center momentum $\mathbf q_B=(\pi/2,\pi/2,\pi/2)$ the massive step
is a non-scalar involution for every mass angle: explicit nonzero
eigenvectors with eigenvalues $+1$ and $-1$ persist for all $\theta$.  The
ordered walk therefore retains exact quasienergies $0$ and $\pi$ at every
mass angle and admits no uniform Floquet gap.
\end{theorem}
\noindent\Kernel{}\quad By Theorem~\ref{thm:class}, on the massive
principal branch these body-center-type points are the \emph{only}
crossings; the obstruction is thus exactly located, not merely exhibited.

\begin{corollary}[Massless corner aliases]\label{cor:aliases}
At $\theta=0$ the eight corners $\mathbf q\in\{0,\pi\}^3$ satisfy
$U(\mathbf q,0)=(-1)^{r}\id$ with $r$ the number of $\pi$ entries; the
three even-parity corners $(\pi,\pi,0)$, $(\pi,0,\pi)$, $(0,\pi,\pi)$ carry
exactly the origin's zero-quasienergy symbol.
\end{corollary}

\section{Which repairs are impossible: two no-go theorems and a corollary}
\label{sec:nogo}

The most economical proposed repair adds a stay-put channel to each
nearest-neighbor factor.

\begin{theorem}[Degree-one stationary no-go]\label{thm:stationary}
Let $F(q)=e^{\ii q}A+B+e^{-\ii q}C$ be unitary for every real $q$ with
$F(0)=\id$ and $F'(0)=-\ii M$ for a Hermitian involution $M$.  Then $B=0$.
The specializations $M=\alpha_j$ for each live axis generator are checked
separately, and a partial-channel control shows the hypothesis is sharp: a
tangent with a stationary kernel admits an exactly unitary degree-one
witness with $B\neq0$.
\end{theorem}
\noindent\Kernel{}

\begin{theorem}[All-coins alias persistence]\label{thm:allcoins}
In the four-channel, range-one, single-factor-per-axis architecture ---
each axis factor unitary at every momentum and carrying the exact
involutory Dirac tangent --- with origin normalization, every
momentum-independent onsite coin $Q$ leaves all
three even-parity corner aliases intact:
$U_Q(0,0,0)=U_Q(\pi,\pi,0)=U_Q(\pi,0,\pi)=U_Q(0,\pi,\pi)$.
\end{theorem}
\noindent\Kernel{}\quad Changing the onsite mass data cannot repair the
cubic partners inside the minimal architecture; extra range, coupled
substeps, enlarged or ancillary cells, or a changed tangent are structurally
required.

\begin{corollary}[Constraint on stay-put doubler-free families]
\label{cor:gs}
The doubler- and pseudo-doubler-free nearest-neighbor family of Gupta and
Short \cite{GuptaShortDoubling} is exactly of the degree-one Laurent form
covered by Theorem~\ref{thm:stationary}: one-step, unitary at every
momentum, normalized at the origin, with stationary block
$\gamma_0=\id-\gamma_+-\gamma_-\neq0$ built from tilted velocity projectors
(their Eqs.~(29)--(30), (37)).  Its tangent generator, though automatically
Hermitian, therefore cannot be an involution ($M^2\neq\id$): stay-put
freedom is purchased exactly by relinquishing the involutory unit-speed
Dirac tangent.  Their $3+1$ family moreover retains extraneous Weyl-like
low-energy solutions at isolated momenta $\pm\mathbf q(\theta)$ (their
Appendix~F), so a strictly local, exactly unitary, fully alias-free $3+1$
walk with the complete involutory Dirac tangent remains open in both
programs.
\end{corollary}

The three results triangulate the design space: inside the minimal
architecture the aliases are unavoidable (Theorem~\ref{thm:allcoins}); a
stay-put escape is incompatible with the exact unit-speed tangent
(Theorem~\ref{thm:stationary}, Corollary~\ref{cor:gs}); and the massive
spectrum of the minimal walk fails to gap at exactly classified points
(Theorems \ref{thm:det}--\ref{thm:bodycenter}).  A Wilson-type Hamiltonian
comparison removes the non-origin corners and opens a uniform massive gap
\emph{at the Hamiltonian level only}; its finite-time exponential is not a
strictly finite-range one-step update, and we do not conflate the two.

\section{Scope, wagers, and what this letter is not}
\label{sec:scope}

These are boundary theorems, not a doubler-free construction.  Two further
kernel-checked boundaries are part of the same suite: the natural positive
homogeneous action of the underlying null-spinor data selects only the
collapsed zero scale (no dynamical mass-scale selection), and exact
two-step temporal blocking closes the mass-only subgroup but not the
ordered split family.  A companion analysis of walks with
spinor-derived defects addresses the corresponding protection question.  The companion paper \cite{CompanionDerivation}
derives the walk's complex mass coin from null-spinor geometry and proves
the corresponding $1+1$ and compact-support continuum control; no result is
headlined in both documents.

The letter's falsifiable wager is Theorem~\ref{thm:allcoins}'s reach: no
four-channel, dimension-two-per-factor, range-one, single-substep
factorized walk achieves alias-free $3+1$ Dirac dynamics; the first
alias-free strict walk must grow a named resource (range, cell, substeps,
or tangent).  Exhibiting an alias-free walk inside the minimal class would
refute the wager and the architecture reading of
Theorem~\ref{thm:allcoins}; eliminating the Gupta--Short residual states
without growing a resource would refute the triangulation.  Equations
\eqref{eq:detfact} and Theorem~\ref{thm:bodycenter} are the regression
tests any proposed successor must pass.

\appendix

\section{Machine verification and artifact}

The formal development uses Lean~4.28.0 and Mathlib \cite{Lean4,Mathlib}.
Every theorem above is a kernel-checked declaration with a build-enforced
assumption footprint contained in
$\{\mathtt{propext},\mathtt{Classical.choice},\mathtt{Quot.sound}\}$.

\begin{center}
\small
\begin{tabularx}{\linewidth}{@{}p{0.42\linewidth}X@{}}
\toprule
Declaration & Statement\\
\midrule
\texttt{det\_splitStep\_sub\_one}, \texttt{det\_splitStep\_add\_one}
 & the factorizations \eqref{eq:detfact}, all momenta and mass angles\\
\texttt{FullBlochZeroClassification} (both iff theorems)
 & Theorem~\ref{thm:class} with $\cos\theta$ boundary controls\\
\texttt{body\_center\_walk\_sq},
\texttt{body\_center\_plus/minus\_mode\_eigen}
 & Theorem~\ref{thm:bodycenter} eigenmodes, every mass angle\\
\texttt{massless\_corner\_parity\_classification},
\texttt{origin\_alias\_*}
 & Corollary~\ref{cor:aliases} with three distinctness witnesses\\
\texttt{stationary\_forces\_zero}, \texttt{dirac\_no\_go},
\texttt{alpha1/2/3\_stationary\_forces\_zero}
 & Theorem~\ref{thm:stationary} and its live-axis specializations\\
\texttt{relaxed\_witness}
 & sharpness control for Theorem~\ref{thm:stationary}\\
\texttt{live\_degree\_one\_factorized\_lower\_bound}
 & Theorem~\ref{thm:allcoins}\\
\texttt{no\_constant\_onsite\_equivalence\_of\_corner\_separation}
 & onsite conjugacy cannot erase a corner alias\\
\texttt{WilsonDiracRegulator.H\_sq},
\texttt{massless\_energy\_eq\_zero\_iff}
 & Hamiltonian-scope Wilson comparison\\
\texttt{no\_nonzero\_stationary\_scale},
\texttt{massCoin4\_block2\_closes},
\texttt{xy\_blocking\_nonclosure\_control}
 & scope boundaries of Section~\ref{sec:scope}\\
\bottomrule
\end{tabularx}
\end{center}

A pre-registered high-momentum benchmark (fixed inputs, $10^{-10}$ kill
threshold) validates the implementation of the theorem chain numerically;
the Lean identities, not the benchmark, carry the proof.  The archival
release --- complete buildable module set, tagged commit, guard outputs,
and deterministic benchmark manifest --- is a stated release gate shared
with the companion paper.

\begin{thebibliography}{9}

\bibitem{NielsenNinomiya}
H.~B.~Nielsen and M.~Ninomiya,
\emph{Absence of neutrinos on a lattice. (I) Proof by homotopy theory},
Nucl. Phys. B 185 (1981) 20--40.

\bibitem{DArianoPerinotti}
G.~M.~D'Ariano and P.~Perinotti,
\emph{Derivation of the Dirac equation from principles of information
processing}, Phys. Rev. A 90 (2014) 062106, arXiv:1306.1934.

\bibitem{MlodinowBrun}
L.~Mlodinow and T.~A.~Brun,
\emph{Discrete spacetime, quantum walks, and relativistic wave equations},
Phys. Rev. A 97 (2018) 042131, arXiv:1802.03910.

\bibitem{MlodinowBrunNoGo}
L.~Mlodinow and T.~A.~Brun,
\emph{Quantum field theory from a quantum cellular automaton in one spatial
dimension and a no-go theorem in higher dimensions},
Phys. Rev. A 102 (2020) 042211, arXiv:2006.08927.

\bibitem{GuptaShortDoubling}
C.~Gupta and A.~J.~Short,
\emph{Fermion doubling in Dirac quantum walks},
Phys. Rev. A, accepted (2026), arXiv:2601.15885, doi:10.1103/kkz3-fthm.

\bibitem{GuptaShortQCA}
C.~Gupta and A.~J.~Short,
\emph{Fermion doubling in quantum cellular automata}, arXiv:2505.07900.

\bibitem{NzonganiTetra}
U.~Nzongani, N.~Eon, I.~Marquez-Martin, A.~Perez, G.~Di Molfetta, and
P.~Arrighi,
\emph{Dirac quantum walk on tetrahedra},
Phys. Rev. A 110 (2024) 042418, arXiv:2404.09840.

\bibitem{FreeQCA}
A.~Bisio, G.~M.~D'Ariano, P.~Perinotti, and A.~Tosini,
\emph{Free quantum field theory from quantum cellular automata},
Found. Phys. 46 (2016) 1032--1081, arXiv:1601.04832.

\bibitem{CompanionDerivation}
Null-Edge Formalization Project,
\emph{Null-spinor area as the rest gap of an exactly unitary Dirac walk},
companion paper, in preparation (2026).

\bibitem{Lean4}
L.~de Moura and S.~Ullrich,
\emph{The Lean 4 theorem prover and programming language},
CADE-28, LNCS 12699 (2021) 625--635.

\bibitem{Mathlib}
The Mathlib Community,
\emph{The Lean mathematical library}, CPP 2020, 367--381.

\end{thebibliography}

\end{document}
